\chapter{Project Management and Development Evidence}
\label{ch:evidence}

\section{Work Completed}
\begin{enumerate}
    \item Android capture app with ARCore pose/depth/quality logging, live IQA, coverage overlay, upload UI.
    \item Reconstruction pipeline: IMU fallback, gyro-aware frame selection, COLMAP SfM with Sim3 metric alignment, Nerfstudio splatfacto training, model export.
    \item FastAPI processing server with local and RunPod cloud GPU modes.
    \item RunPod integration with guaranteed pod lifecycle management.
    \item CI/CD pipeline (GitHub Actions): lint, unit tests, APK build, web build, hygiene checks.
    \item Extensive root-cause analysis of degenerate splat geometry (6 causes, documented in \texttt{Training\_Analysis.md}).
    \item Server launcher script (\texttt{start\_server.bat}) with \texttt{.env} support.
\end{enumerate}

\section{Work Remaining}
\begin{enumerate}
    \item \textbf{Web viewer}: Three.js/WebGPU Gaussian splat viewer with orbit controls and measurement tools.
    \item \textbf{Measurement tool}: Point-to-point distance measurement in the metric 3D model.
    \item \textbf{Structural layer}: Complete \texttt{structure.py}: wall-to-wall dimensions, doors and openings, ceiling-height detection, and automatic floor-plan export.
    \item \textbf{Wall snapping}: Measurement clicks that snap to structural-layer walls, and a floor-plan view in the web viewer.
    \item \textbf{3DGS refinement}: One-command v6 pipeline, validation on the second room, and IMU-driven motion-blur modelling.
    \item \textbf{3D mesh creation}: In scope as planned work; not yet scheduled.
    \item \textbf{Field validation}: Metric accuracy validation on $\geq$3 additional rooms with varied geometry.
    \item \textbf{User testing}: Usability evaluation with external users (non-team members).
\end{enumerate}

\section{Individual Responsibilities}

The primary evidence of each member's contribution is the set of modules they own: the files whose commit history carries their name (\texttt{git log -- <file>}). Table~\ref{tab:responsibilities} lists them.

\begin{table}[h]
\centering
\caption{Owned modules per team member (from the commit history of each file)}
\label{tab:responsibilities}
\begin{tabular}{p{3.2cm}p{9.3cm}}
\toprule
\textbf{Team Member} & \textbf{Owned modules} \\
\midrule
Muhammad Umair (23I-0662) & Processing server: \texttt{server/main.py}, \texttt{pipeline.py}, \texttt{runpod\_client.py}, \texttt{remote\_gpu.py}. Android upload and app screens: \texttt{ProcessingActivity}, \texttt{ScanDetailActivity}, \texttt{UpdateChecker}; co-author of \texttt{HelloArActivity}, \texttt{CoverageMap}, \texttt{CoverageOverlayView}, \texttt{QualityMonitor}. Web UI prototype (\texttt{prototype/}). GPU Docker workflow; co-author of the CD workflow. Server and job scripts (\texttt{scripts/}). Proposal, report and project documentation. \\
\midrule
Zaid Fareed (23I-0829) & CI/CD: \texttt{ci.yml}, \texttt{cd.yml}, APK/manifest checker (\texttt{.github/android/}), \texttt{ruff.toml}, \texttt{pytest.ini}, pull request template. The whole pytest suite (\texttt{tests/}, all 12 files and \texttt{conftest.py}). Trainer experiments and fixes (\texttt{train.py} checkpoints and depth loss, \texttt{mvs.py}), ARKitScenes trainer experiments (\texttt{docs/Experiments\_ARKitScenes.md}). CI/CD documentation. \\
\midrule
Muhammad Tayyab (23I-0636) & Reconstruction pipeline (\texttt{src/arcore3dgs/}): \texttt{fallback.py}, \texttt{imu.py}, \texttt{relocalize\_pnp.py}, \texttt{prepare.py}, \texttt{select\_sharp.py}, \texttt{ba.py}, \texttt{sfm\_colmap.py}, \texttt{export\_nerfstudio.py}, \texttt{mvs\_init.py}, \texttt{structure.py}, evaluation tools (\texttt{fair\_compare.py}, \texttt{fit\_diag.py}, \texttt{diagnose.py}). Capture logging and live IQA in the app: \texttt{CaptureLogger}, \texttt{QualityMonitor}; co-author of \texttt{HelloArActivity} and the coverage overlay. Training analysis (\texttt{Training\_Analysis.md}), tape-measure validation and second-room evaluation. \\
\bottomrule
\end{tabular}
\end{table}

\section{Contribution of Each Team Member}

\paragraph{Muhammad Umair (23I-0662):}
Project lead. Built the processing server (FastAPI endpoints, job model with per-step timeouts, local pipeline runner) and the RunPod integration (pod creation, SSH readiness, remote pipeline, termination in \texttt{finally} blocks). Built the app's upload and progress screen and the scan gallery/detail screens, worked on the main AR activity and coverage overlay, built the web UI prototype, and wrote the proposal and report.

\paragraph{Zaid Fareed (23I-0829):}
Built the CI/CD infrastructure (lint, test and build workflows on three Python versions, Android build and signing checks, repository hygiene, releases) and wrote the entire automated test suite that backs the ``Pass (CI)'' results in Chapter~\ref{ch:testing}. Also contributed trainer fixes and the ARKitScenes trainer experiments.

\paragraph{Muhammad Tayyab (23I-0636):}
Built the reconstruction pipeline: IMU fallback and re-alignment, keyframe and gyro-aware frame selection, bundle adjustment, the COLMAP SfM stage with Sim3 metric alignment, dense MVS initialisation, the structural layer, and the evaluation tools. Wrote the app's capture logging and live IQA, ran the training analysis, and produced the v6 model, the tape-measure validation and the second-room result.

\paragraph{Commit counts.}
The repository has \textbf{133} commits: Muhammad Umair 85, Zaid Fareed 26, Muhammad Tayyab 20 (committing as ``BlazingBeam''), and 2 automated commits by the GitHub Actions bot. \textbf{These counts reflect differing commit granularity, not differing effort.} Umair commits each small change separately; Tayyab works in large batches, so a single commit of his often carries a whole pipeline stage (for example, the commit that introduced image-based SfM poses and lens self-calibration, or the one that introduced gyro-aware frames, dense MVS initialisation and the structural layer). The owned modules above, not the commit counts, are the evidence of contribution.

\section{Repository and Version Control Evidence}
The project repository is hosted on GitHub. Key statistics:
\begin{itemize}
    \item \textbf{Repository}: \texttt{github.com/Umairishuman/FYP} (private)
    \item \textbf{Branches}: \texttt{master} (stable) and work branches such as \texttt{ci/github-actions}, \texttt{ci/pipeline-updates} and \texttt{exp/trainer-arkitscenes}
    \item \textbf{Commit history}: 133 commits (85 / 26 / 20, plus 2 bot commits)
    \item \textbf{CI}: GitHub Actions runs lint, tests and builds on work branches and pull requests
    \item \textbf{Code review}: Pull request template with area checkboxes and a testing checklist
\end{itemize}

\section{Milestones and Commit Evidence}

\begin{table}[h]
\centering
\caption{Key milestones with commit evidence}
\label{tab:milestones}
\begin{tabular}{p{2.5cm}p{5cm}p{4cm}}
\toprule
\textbf{Milestone} & \textbf{Description} & \textbf{Evidence} \\
\midrule
M1: Capture & ARCore recording with pose/depth/IMU logging & Commits by Tayyab, Umair \\
M2: Pipeline & IMU fallback, COLMAP SfM, splatfacto training & Commits by Tayyab; trainer fixes by Zaid \\
M3: Quality & Live IQA, coverage overlay, gyro-aware selection & IQA and selection by Tayyab; coverage overlay by Umair and Tayyab \\
M4: Server & FastAPI + RunPod integration, upload UI & Commits by Umair \\
M5: CI/CD & GitHub Actions workflows, test suite & Commits by Zaid; workflow merges by Umair \\
M6: Validation & v3--v6 results, tape measure, second room & Commits by Tayyab; numbers in \texttt{docs/progress/} \\
\bottomrule
\end{tabular}
\end{table}

\section{Tools and Tool-Assisted Development}

\textbf{IMPORTANT: Generative AI Tool Disclosure}

In compliance with university policy, we disclose the use of the following AI-assisted development tools:

\begin{table}[h]
\centering
\caption{AI and automated tools used during development}
\label{tab:aitools}
\begin{tabular}{p{3cm}p{4cm}p{5cm}}
\toprule
\textbf{Tool} & \textbf{Usage} & \textbf{Scope / Boundaries} \\
\midrule
Claude Code (Anthropic) & Agentic coding assistant, used in local and cloud sessions & Code generation and refactoring, debugging and error diagnosis, CI/CD workflow changes (some commits carry a \texttt{Co-Authored-By: Claude} trailer), project documentation, drafting of this report, and the correction pass on this report after supervisor review (checking every figure against \texttt{docs/progress/} and the team's decisions, rewriting the test and meeting sections, and drawing the TikZ diagrams). All output was reviewed, tested and edited by team members, who remain responsible for it. \\
GitHub Copilot & Inline code completion & Used for auto-completion in Android Studio and VS Code. All suggestions were reviewed before acceptance. \\
ChatGPT (OpenAI) & Research queries, algorithm explanations & Used to understand COLMAP parameters, quaternion math, and coordinate system conventions. No code directly copied. \\
\bottomrule
\end{tabular}
\end{table}

\paragraph{Statement of originality:}
All core algorithms (Sim3 alignment, gyro-aware frame selection, coverage map, IQA pipeline, RunPod lifecycle management) were designed by the team. AI tools were used as assistants for code generation, debugging, explanation and report drafting, as disclosed in Table~\ref{tab:aitools}. The system architecture, experimental methodology, root-cause analysis, and design decisions are entirely our own work. All AI-generated code was reviewed, understood, tested, and adapted by team members before integration into the codebase.

\paragraph{Other development tools:}
\begin{itemize}
    \item Android Studio (Ladybug) with Android Gradle Plugin
    \item Visual Studio Code with Python and Java extensions
    \item Git + GitHub for version control
    \item Postman for API testing during server development
    \item ADB (Android Debug Bridge) for device-level debugging
    \item NVIDIA Nsight for GPU profiling during training experiments
\end{itemize}

\section{External Components and Dependencies}

\begin{table}[h]
\centering
\caption{External components and their licenses}
\label{tab:dependencies}
\begin{tabular}{p{3.5cm}p{2cm}p{2.5cm}p{3cm}}
\toprule
\textbf{Component} & \textbf{Version} & \textbf{License} & \textbf{Usage} \\
\midrule
ARCore SDK & 1.44+ & Apache 2.0 & Pose tracking, depth \\
COLMAP & 4.2 & BSD-3 & Structure from Motion \\
Nerfstudio & 1.1.5 & Apache 2.0 & 3DGS training \\
FastAPI & 0.110+ & MIT & REST API server \\
OkHttp & 4.12.0 & Apache 2.0 & HTTP client \\
PyTorch & 2.1+ & BSD-3 & GPU computation \\
OpenCV & 4.9+ & Apache 2.0 & Image processing \\
\bottomrule
\end{tabular}
\end{table}

All external components are used under permissive open-source licenses compatible with academic use. No proprietary or restrictively-licensed code is included.

\section{Risks and Mitigation}

\begin{table}[h]
\centering
\caption{Risk register}
\label{tab:risks}
\begin{tabular}{p{3cm}p{1.5cm}p{1.5cm}p{5cm}}
\toprule
\textbf{Risk} & \textbf{Likelihood} & \textbf{Impact} & \textbf{Mitigation} \\
\midrule
Cloud GPU costs exceed budget & Medium & High & Per-step timeouts; guaranteed pod termination in finally block; GPU type selectable per job via \texttt{/config} (default RTX 4090) \\
ARCore API deprecation & Low & High & Abstracted behind CaptureLogger interface; can swap to AREngine \\
COLMAP SfM fails on difficult scenes & Medium & Medium & IMU fallback for tracking gaps; live IQA reduces bad input; manual SfM parameter overrides available \\
Nerfstudio breaking changes & Low & Medium & Pinned to version 1.1.5; tested upgrade path \\
Team member unavailability & Low & Medium & All code in shared repo; documented in Flow.md and TODO.md \\
\bottomrule
\end{tabular}
\end{table}

\section{Security and Privacy Considerations}
\begin{itemize}
    \item \textbf{API key management}: The RunPod API key is stored in a \texttt{.env} file excluded from version control via \texttt{.gitignore}. The key is loaded at runtime from environment variables, never hardcoded.
    \item \textbf{CI secrets scanning}: The GitHub Actions CI pipeline includes a secrets detection check that blocks merges if API keys, tokens, or credentials are found in committed files.
    \item \textbf{Data privacy}: Captured scan data (video, depth, poses) is transmitted to the processing server over the local network. In production, HTTPS would be required. Scan data is stored only for the duration of processing and can be deleted via the API.
    \item \textbf{Cloud GPU security}: RunPod pods are created on-demand, used for a single job, and terminated immediately after. No persistent data remains on the cloud GPU.
\end{itemize}

\section{Current Limitations}
\begin{enumerate}
    \item \textbf{Mirrors and glass}: Highly reflective and transparent surfaces cause SfM outliers and incorrect depth. The model shows artifacts in these regions. Planned FYP-2 investigation: detect and mask reflective regions.
    \item \textbf{Large rooms}: Rooms exceeding approximately 8$\times$8\,m may exceed GPU memory during training on the RTX 3070 (8\,GB). A 24\,GB cloud GPU (RTX 4090, the RunPod default) is expected to handle larger scenes; this has not yet been measured.
    \item \textbf{Outdoor scenes}: The system is designed for indoor environments. Outdoor scenes with sky, vegetation, and long-range views are out of scope.
    \item \textbf{Real-time processing}: Processing takes approximately 75 minutes on the development RTX 3070; no cloud-GPU timing has been measured yet. Real-time reconstruction is not feasible with current 3DGS training times.
    \item \textbf{Single-floor only}: The current coverage map and measurement tools assume a single-floor environment. Multi-floor buildings require manual segmentation.
    \item \textbf{No web viewer yet}: The 3D model can currently only be viewed in a local desktop viewer. The web-based Three.js viewer is planned for FYP-2.
\end{enumerate}
